% ---------------------------------------------------------------------------
% The framework, compressed to what the measurement sections use. Definitions are the
% complete paper's, verbatim in content; the proofs they carry there are cited. The
% two notes this version adds weight to are the ones its result turns on: the kind
% sits on the series rather than in the query, and a caveat set is part of what the
% emitter returns.
% ---------------------------------------------------------------------------

\section{Formal Framework}
\label{s2:sec:framework}

``Are these languages inter-translatable'' is not yet a question. It becomes one once
four things are fixed: what the systems store, what a query denotes on what they store,
when two denotations count as the same, and what a translator is permitted to do when it
cannot proceed.

\begin{problem}
Let $L_s$ and $L_t$ be query languages over time-series databases related by an
ingestion mapping rather than shared. Determine the largest fragment of $L_s$ that admits
a meaning-preserving translation into $L_t$, exhibit the obstructions outside it, and
measure what share of queries written in practice the fragment contains.
\end{problem}

\noindent
Three features of that statement shape everything after it: it asks for a \emph{fragment},
so the answer is a set of queries and not a yes or no; it asks for the obstructions, so a
translator that fails must say why in a form that can be counted; and it asks what share of
\emph{practice} the fragment covers, which makes the corpus part of the result rather than
an illustration of it.

\subsection{What the systems store}

\begin{definition}[Series, database, result]\label{s2:def:model}
A \emph{series} is a pair $s=(\lambda,\langle (t_1,v_1),\dots,(t_n,v_n)\rangle)$ of a
finite label set $\lambda$ --- a partial map from label names to strings --- and samples
at strictly increasing times $t_i\in\mathbb{T}$ with values $v_i\in\mathbb{R}$. A series
carries a \emph{kind}: $\mathsf{counter}$ for a cumulative total, $\mathsf{delta}$ for a
per-interval increment, $\mathsf{gauge}$ otherwise. A \emph{database} $D$ is a finite set
of series with pairwise distinct label sets, and $\mathcal{D}$ is the class of all
databases. A \emph{result} is a finite partial map from label sets to $\mathbb{R}$.
\end{definition}

\noindent
Two choices here are not neutral, and the second is this paper's subject. Requiring distinct
label sets makes a label set the identity of a series, which is what both systems assume and
what lets a result be keyed by one; and the kind is a property of the \emph{series} rather
than of the query --- no query in either language declares it, the translation depends on
it, and \S\ref{s2:sec:unsaid} measures what it costs to guess.

\subsection{Queries, denotations, translators}

A \emph{query language} is a pair $(G,\sem{\cdot})$ of a context-free grammar $G$ and a
denotation $\sem{\cdot}:\Lang(G)\times\mathcal{D}\times\mathbb{T}\to\mathit{Result}$,
written $\Lang(G_{\mathrm{DD}})$, $\Lang(G_{\mathrm{PQ}})$ and
$\sem{\cdot}_{\mathrm{DD}}$, $\sem{\cdot}_{\mathrm{PQ}}$ for the two here. Taking the set
of strings to be generated by a grammar is what makes the grammar question one about
formal languages rather than about parsers, and taking the meaning to be a function of a
database and an instant is what lets two queries be compared at all.

One thing that pair declares and does not define, and the paper is better for saying so.
The two denotations are \emph{declared}: no definition of $\sem{\cdot}_{\mathrm{DD}}$ or
$\sem{\cdot}_{\mathrm{PQ}}$ appears anywhere, and the only denotation written out is the
intermediate representation's. So every soundness statement here is an equation between
two functions pinned down by measurement rather than by definition --- the complete paper
puts Datadog's bucket boundaries, its reset handling and the way it pairs operands to a
live account, and takes PromQL's side from the published specification. A measured
semantics is the honest source, but the theorems are conditional on those measurements.

\begin{definition}[Translator]\label{s2:def:translator}
A \emph{translator} is a total function $T:\Lang(G_s)\to\Lang(G_t)\cup\{\bot\}$, where
$\bot$ carries a \emph{gap code} naming the construct that stopped it. Its \emph{domain}
is $\mathrm{dom}(T)=\{q : T(q)\neq\bot\}$, the fragment of the Problem Statement. The
translator this paper builds factors through an intermediate representation, so it is
useful to name the halves: a \emph{lowering} $L:\Lang(G_s)\to\mathit{IR}\cup\{\bot\}$ and
an \emph{emission} $E:\mathit{IR}\to(\Lang(G_t)\times\mathcal{P}(\text{codes}))
\cup\{\bot\}$, with $T=\pi_1\circ E\circ L$ and $\bot$ absorbing.
\end{definition}

\noindent
Totality is one substantive requirement and the second component of $E$ is the other. A
translator that may crash, or silently emit something plausible, cannot be measured.
Requiring every failure to be a $\bot$ with a code makes $\mathrm{dom}(T)$ observable and
makes the complement a \emph{catalogue} rather than a list of bug reports; a query is
abandoned at its first obstruction, so the codes count first refusals. The set of codes $E$
returns \emph{beside} a successful expression is the other half of the instrument, and the
one this paper's result comes out of: it records what the emitter had to assume, so an
assumption is counted rather than conceded. Every figure in \S\ref{s2:sec:unsaid} is a
share of that set.

\begin{remark}[A refusal is not evidence by itself]\label{s2:rem:refusal}
A gap code looks like a statement about the languages, and it is in the first instance a
statement about $T$: the construct stopped \emph{this} translator. The two coincide only
when something independent says so --- a proof that no target expression exists, a decision
procedure that agrees, or an implementation that removes the refusal and so shows it was
never a gap. \S\ref{s2:sec:coverage} reports eight occasions on which this paper had the two
confused, every one in the direction that would have made its thesis look better.
\end{remark}

\subsection{When two results agree, and what soundness means}

Comparing results needs a distance, not an equality test: the two systems compute in
floating point over different code paths, so bit equality would report disagreement where
none is meant, and relative error is the right scale because the quantities span many
orders of magnitude. We write $d(r_1,r_2)$ for the complete paper's result distance:
$\infty$ when the two results are keyed by different label sets, and otherwise the greatest
relative error $|a-b|/\max(|a|,|b|)$ over the keys they share, with two empty results at
distance $0$. Differing label sets give $\infty$ rather than a large number because the
translation returned the wrong \emph{series}, and no tolerance should absorb that.

\begin{definition}[Exact and $\varepsilon$-sound]\label{s2:def:sound}
Let $\mathcal{C}\subseteq\mathcal{D}$ and let $u:\mathbb{T}\to\mathbb{T}$ be an
\emph{evaluation map}. A translator $T$ is \emph{$\varepsilon$-sound on $\mathcal{C}$
along $u$} if for every $q\in\mathrm{dom}(T)$, every $D\in\mathcal{C}$ and every
$t\in\mathbb{T}$,
\[
d\big(\sem{q}_s(\phi_\rho(D),t),\ \sem{T(q)}_t(D,u(t))\big)\le\varepsilon,
\]
and \emph{exact} on $\mathcal{C}$ along $u$ when $\varepsilon=0$.
\end{definition}

\noindent
The $\phi_\rho$ is the ingestion mapping of \S\ref{s2:sec:intro}: the two systems do not
receive the same $D$, so the source query is evaluated on the database the source system
would have been given. Setting $\phi_\rho=\mathrm{id}$ recovers the naive definition, and
the complete paper builds $\phi_\rho$. Everything claimed here is exact on a class
$\mathcal{C}$ the claim states --- a deliberate choice of target, since an
$\varepsilon$-sound translation of a dashboard is a dashboard whose numbers are nearly
right, which is a worse engineering outcome than a refusal because nothing downstream can
tell. And $\mathcal{C}$ is where the assumptions of \S\ref{s2:sec:unsaid} live: a theorem
exact on $\mathcal{C}$ says nothing about a database outside it, and whether the corpus's
databases are inside it is not something the query text settles.

\begin{definition}[Coverage]\label{s2:def:kappa}
Let $Q$ be a finite multiset of queries --- a corpus, with multiplicity, so a query used
on forty dashboards counts forty times. The \emph{coverage} of $T$ on $Q$ is
$\kappa(T,Q)=|\{q\in Q : T(q)\neq\bot\}|/|Q|$.
\end{definition}

\noindent
This counts queries the translator is \emph{willing} to translate, which is the weaker of
two measures: the stronger runs $q$ against the source system and $T(q)$ against the target
on the same data and counts a query only when the two results agree. The complete paper runs
that differential measure against a real Datadog account and a local Prometheus over every
distinct query shape, and the two agree on all $330$ of the $333$ shapes it could compare.
Every $\kappa$ quoted here is Definition~\ref{s2:def:kappa}, and the gap between the two
measures is the subject of \S\ref{s2:sec:unsaid}.
